% GENERATED FILE. Edit canonical lessons, then run npm run generate:books.
% canonical-source-hash: fnv1a64:f9c5afff03ab7c89
% canonical-lessons: KA-C271-pincani, KA-C271-danda, KA-C271-muskara, KA-C271-samsthe, KA-C271-ilakhe, KA-R271-first-pass-words-from-chapters-267-268, KA-R271-second-pass-words-from-chapters-269-271

\chapter{Pension, Fine, Strike, Organisation, Government department}
\label{ch:ka-pension-fine-strike-organisation-government-department}
\hlchaptermodality{271}

\begin{chapteropening}
\textbf{By the end of this chapter:} \emph{I can say a pension, a fine, a strike, an organisation and a government department.}

Complete the last lesson of chapter 271: I can say a pension, a fine, a strike, an organisation and a government department.
\end{chapteropening}

\section[\texorpdfstring{\kn{ಪಿಂಚಣಿ} (piñcaṇi)}{ಪಿಂಚಣಿ (piñcaṇi)}]{\kn{ಪಿಂಚಣಿ} (piñcaṇi) — a pension}

\label{lesson:KA-C271-pincani}



\begin{quote}
Before the new one: say the Kannada for a barber, then the Kannada for a washerman.
\end{quote}

\subsection*{You'll want to know: \kn{ಪಿಂಚಣಿ}}
\textbf{\kn{ಪಿಂಚಣಿ}} — \emph{piñcaṇi} — \textquotedblleft{}a pension\textquotedblright{}.

\begin{grammarlens}[title={What you've built}]
One more word for this chapter.
\end{grammarlens}

\subsection*{Your turn}
\begin{itemize}
\raggedright
  \item \emph{Say it:} \emph{piñcaṇi}
  \item \emph{Say it:} \emph{piñcaṇi}, once more
  \item \emph{Say it:} say \emph{piñcaṇi}
  \item \emph{Recall:} say the Kannada for a barber, then the Kannada for a washerman, then say \emph{piñcaṇi} again
\end{itemize}

\begin{tcolorbox}[breakable,colback=teal!4,colframe=teal!35!black,title={Before you move on}]
What does \kn{ಪಿಂಚಣಿ} mean? (\textquotedblleft{}A pension\textquotedblright{}.) Say it once more.
\end{tcolorbox}
\section[\texorpdfstring{\kn{ದಂಡ} (daṇḍa)}{ದಂಡ (daṇḍa)}]{\kn{ದಂಡ} (daṇḍa) — a fine, a penalty}

\label{lesson:KA-C271-danda}



\begin{quote}
Before the new one: say the Kannada for a washerman, then the Kannada for a pension.
\end{quote}

\subsection*{You'll want to know: \kn{ದಂಡ}}
\textbf{\kn{ದಂಡ}} — \emph{daṇḍa} — \textquotedblleft{}a fine, a penalty\textquotedblright{}.

\begin{grammarlens}[title={What you've built}]
One more word for this chapter.
\end{grammarlens}

\subsection*{Your turn}
\begin{itemize}
\raggedright
  \item \emph{Say it:} \emph{daṇḍa}
  \item \emph{Say it:} \emph{daṇḍa}, once more
  \item \emph{Say it:} say \emph{daṇḍa}
  \item \emph{Recall:} say the Kannada for a washerman, then the Kannada for a pension, then say \emph{daṇḍa} again
\end{itemize}

\begin{tcolorbox}[breakable,colback=teal!4,colframe=teal!35!black,title={Before you move on}]
What does \kn{ದಂಡ} mean? (\textquotedblleft{}A fine, a penalty\textquotedblright{}.) Say it once more.
\end{tcolorbox}
\section[\texorpdfstring{\kn{ಮುಷ್ಕರ} (muṣkara)}{ಮುಷ್ಕರ (muṣkara)}]{\kn{ಮುಷ್ಕರ} (muṣkara) — a strike (by workers)}

\label{lesson:KA-C271-muskara}



\begin{quote}
Before the new one: say the Kannada for a pension, then the Kannada for a fine.
\end{quote}

\subsection*{You'll want to know: \kn{ಮುಷ್ಕರ}}
\textbf{\kn{ಮುಷ್ಕರ}} — \emph{muṣkara} — \textquotedblleft{}a strike (by workers)\textquotedblright{}.

\begin{grammarlens}[title={What you've built}]
One more word for this chapter.
\end{grammarlens}

\subsection*{Your turn}
\begin{itemize}
\raggedright
  \item \emph{Say it:} \emph{muṣkara}
  \item \emph{Say it:} \emph{muṣkara}, once more
  \item \emph{Say it:} say \emph{muṣkara}
  \item \emph{Recall:} say the Kannada for a pension, then the Kannada for a fine, then say \emph{muṣkara} again
\end{itemize}

\begin{tcolorbox}[breakable,colback=teal!4,colframe=teal!35!black,title={Before you move on}]
What does \kn{ಮುಷ್ಕರ} mean? (\textquotedblleft{}A strike (by workers)\textquotedblright{}.) Say it once more.
\end{tcolorbox}
\section[\texorpdfstring{\kn{ಸಂಸ್ಥೆ} (saṃsthe)}{ಸಂಸ್ಥೆ (saṃsthe)}]{\kn{ಸಂಸ್ಥೆ} (saṃsthe) — an organisation, an institution}

\label{lesson:KA-C271-samsthe}



\begin{quote}
Before the new one: say the Kannada for a fine, then the Kannada for a strike.
\end{quote}

\subsection*{You'll want to know: \kn{ಸಂಸ್ಥೆ}}
\textbf{\kn{ಸಂಸ್ಥೆ}} — \emph{saṃsthe} — \textquotedblleft{}an organisation, an institution\textquotedblright{}.

\begin{grammarlens}[title={What you've built}]
One more word for this chapter.
\end{grammarlens}

\subsection*{Your turn}
\begin{itemize}
\raggedright
  \item \emph{Say it:} \emph{saṃsthe}
  \item \emph{Say it:} \emph{saṃsthe}, once more
  \item \emph{Say it:} say \emph{saṃsthe}
  \item \emph{Recall:} say the Kannada for a fine, then the Kannada for a strike, then say \emph{saṃsthe} again
\end{itemize}

\begin{tcolorbox}[breakable,colback=teal!4,colframe=teal!35!black,title={Before you move on}]
What does \kn{ಸಂಸ್ಥೆ} mean? (\textquotedblleft{}An organisation, an institution\textquotedblright{}.) Say it once more.
\end{tcolorbox}
\section[\texorpdfstring{\kn{ಇಲಾಖೆ} (ilākhe)}{ಇಲಾಖೆ (ilākhe)}]{\kn{ಇಲಾಖೆ} (ilākhe) — a government department}

\label{lesson:KA-C271-ilakhe}



\begin{quote}
Before the new one: say the Kannada for a strike, then the Kannada for an organisation.
\end{quote}

\subsection*{You'll want to know: \kn{ಇಲಾಖೆ}}
\textbf{\kn{ಇಲಾಖೆ}} — \emph{ilākhe} — \textquotedblleft{}a government department\textquotedblright{}.

\begin{grammarlens}[title={What you've built}]
One more word for this chapter.
\end{grammarlens}

\subsection*{Your turn}
\begin{itemize}
\raggedright
  \item \emph{Say it:} \emph{ilākhe}
  \item \emph{Say it:} \emph{ilākhe}, once more
  \item \emph{Say it:} say \emph{ilākhe}
  \item \emph{Recall:} say the Kannada for a strike, then the Kannada for an organisation, then say \emph{ilākhe} again
\end{itemize}

\begin{tcolorbox}[breakable,colback=teal!4,colframe=teal!35!black,title={Before you move on}]
What does \kn{ಇಲಾಖೆ} mean? (\textquotedblleft{}A government department\textquotedblright{}.) Say it once more.
\end{tcolorbox}
\section[(dialogue)]{Words from chapters 267-268 — a first pass}

\label{lesson:KA-R271-first-pass-words-from-chapters-267-268}



\begin{quote}
Say the words for an organisation and a government department.
\end{quote}

\subsection*{Your turn}
Say these aloud:

\begin{itemize}
\raggedright
  \item manure, a flowerpot, a pipe, a pump, a gate
  \item a lawyer, an engineer, a postman, a company, a class
\end{itemize}

\begin{tcolorbox}[breakable,colback=teal!4,colframe=teal!35!black,title={Before you move on}]
Manure? (\textbf{\kn{ಗೊಬ್ಬರ}}, \emph{gobbara}.) A flowerpot? (\textbf{\kn{ಕುಂಡ}}, \emph{kuṇḍa}.) A pipe? (\textbf{\kn{ಪೈಪ್}}, \emph{paip}.) A pump? (\textbf{\kn{ಪಂಪು}}, \emph{pampu}.) A gate? (\textbf{\kn{ಗೇಟು}}, \emph{gēṭu}.) A lawyer? (\textbf{\kn{ವಕೀಲ}}, \emph{vakīla}.) An engineer? (\textbf{\kn{ಎಂಜಿನಿಯರ್}}, \emph{eñjiniyar}.) A postman? (\textbf{\kn{ಅಂಚೆಯವನು}}, \emph{añceyavanu}.) A company? (\textbf{\kn{ಕಂಪನಿ}}, \emph{kampani}.) A class? (\textbf{\kn{ತರಗತಿ}}, \emph{taragati}.)
\end{tcolorbox}
\section[(dialogue)]{Words from chapters 269-271 — a second pass}

\label{lesson:KA-R271-second-pass-words-from-chapters-269-271}



\begin{quote}
Say the words for a report and a problem.
\end{quote}

\subsection*{Your turn}
Say these aloud:

\begin{itemize}
\raggedright
  \item a report, a problem, news, information, a mark
  \item geography, a word, a sentence, a barber, a washerman
  \item a pension, a fine, a strike, an organisation, a government department
\end{itemize}

\begin{tcolorbox}[breakable,colback=teal!4,colframe=teal!35!black,title={Before you move on}]
A report? (\textbf{\kn{ವರದಿ}}, \emph{varadi}.) A problem? (\textbf{\kn{ಸಮಸ್ಯೆ}}, \emph{samasye}.) News? (\textbf{\kn{ಸುದ್ದಿ}}, \emph{suddi}.) Information? (\textbf{\kn{ಮಾಹಿತಿ}}, \emph{māhiti}.) A mark? (\textbf{\kn{ಅಂಕ}}, \emph{aṅka}.) Geography? (\textbf{\kn{ಭೂಗೋಳ}}, \emph{bhūgōḷa}.) A word? (\textbf{\kn{ಪದ}}, \emph{pada}.) A sentence? (\textbf{\kn{ವಾಕ್ಯ}}, \emph{vākya}.) A barber? (\textbf{\kn{ಕ್ಷೌರಿಕ}}, \emph{kṣaurika}.) A washerman? (\textbf{\kn{ಧೋಬಿ}}, \emph{dhōbi}.) A pension? (\textbf{\kn{ಪಿಂಚಣಿ}}, \emph{piñcaṇi}.) A fine? (\textbf{\kn{ದಂಡ}}, \emph{daṇḍa}.) A strike? (\textbf{\kn{ಮುಷ್ಕರ}}, \emph{muṣkara}.) An organisation? (\textbf{\kn{ಸಂಸ್ಥೆ}}, \emph{saṃsthe}.) A government department? (\textbf{\kn{ಇಲಾಖೆ}}, \emph{ilākhe}.)
\end{tcolorbox}
